\documentclass{amsart}
\newenvironment{dummy}{}{}
\newenvironment{test}{}{}
\begin{document}
\begin{dummy}
Without the \texttt{--extra\_mark\_envs} option specified, {\tt corrinline} could not run successfully on this file.

So if we have an **ERROUR** in this document, and we annotate it in the PDF, we won't be able to do anything because no tex boxes
will be created by the program since all the text is inside an environment which was not specified for marking.

Most of the time this isn't an issue because the most frequently used environments which are safe to mark in are already compiled
by the program.

Ideally you should never need to use this option, but if for some reason, even after cleanup, there remains
peculiar environments which you know can be marked but are not standard, you can add them to this option like
\begin{verbatim}
corrinline annotated_file.pdf latex_file.tex --extra-marked-envs=test,dummy
\end{verbatim}
\end{dummy}
\begin{test}
You can also do 
\begin{verbatim}
corrinline annotated_file.pdf latex_file.tex --extra-marked-envs="test,dummy"
\end{verbatim}
\end{test}
\end{document}
